\documentclass{beamer}

\usepackage{amsmath,amsfonts,amssymb}
\usepackage{physics}
\usepackage{bm}
\usepackage{quantikz}

\title{Circuit Depth Scaling in Variational Quantum Algorithms}
\subtitle{Trotterization and the Unitary Coupled Cluster Ansatz}
\author{}
\date{}

\begin{document}

\frame{\titlepage}

%------------------------------------------------
\begin{frame}{Outline}

\begin{itemize}
\item Variational quantum eigensolvers
\item Exponentials of operator sums
\item Suzuki–Trotter decompositions
\item Baker–Campbell–Hausdorff error analysis
\item Circuit depth scaling
\item Application to unitary coupled cluster
\item Adaptive ansätze
\item Parametric matrix models
\end{itemize}

\end{frame}

%================================================
\section{Variational Quantum Eigensolver}

%------------------------------------------------
\begin{frame}{Variational Quantum Algorithms}

The variational quantum eigensolver (VQE) prepares states

\[
|\Psi(\theta)\rangle = U(\theta)|\Phi_0\rangle
\]

minimizes

\[
E(\theta)=\langle\Psi(\theta)|H|\Psi(\theta)\rangle .
\]

\end{frame}

%------------------------------------------------
\begin{frame}{Ansatz Circuits}

  Common ansätze include

  \begin{itemize}
  \item hardware efficient circuits
  \item time evolution inspired circuits
  \item unitary coupled cluster
  \end{itemize}

  Many ansätze are generated by exponentials

  \[
  U(\theta)=e^{-iA}
  \]

\end{frame}

%================================================
\section{Exponentials of Operator Sums}

%------------------------------------------------
\begin{frame}{Operator Decomposition}

  Suppose

  \[
  A = \sum_k A_k
  \]

  In general

  \[
    [A_i,A_j] \neq 0
    \]

    Therefore

    \[
    e^{A+B} \neq e^A e^B
    \]

\end{frame}

%------------------------------------------------
\begin{frame}{Trotter Approximation}

  We approximate

  \[
  e^{\sum_k A_k}
  \]

  by

  \[
  \left(
  \prod_k e^{A_k\Delta t}
  \right)^n
  \]

  where

  \[
  n\Delta t = 1
  \]

\end{frame}

%------------------------------------------------
\begin{frame}{Number of Trotter Steps}

  Since

  \[
  n\Delta t = 1
  \]

  we obtain

  \[
  n=\frac{1}{\Delta t}
  \]

\end{frame}

%------------------------------------------------
\begin{frame}{Circuit Interpretation}

  Each factor

  \[
  e^{A_k\Delta t}
  \]

  corresponds to a circuit block.

  One Trotter step contains

  \[
  M
  \]

  blocks if

  \[
  A=\sum_{k=1}^M A_k
  \]

\end{frame}

%------------------------------------------------
\begin{frame}{Circuit Depth}

  Total depth

  \[
  D \sim n M
  \]

  Using

  \[
  n = \frac{1}{\Delta t}
  \]

  we obtain

  \[
  D \sim \frac{M}{\Delta t}
  \]

\end{frame}

%================================================
\section{Baker–Campbell–Hausdorff Analysis}

%------------------------------------------------
\begin{frame}{BCH Formula}

  The BCH formula states

  \[
  e^A e^B
  =
  e^{A+B+\frac12[A,B]+\frac1{12}[A,[A,B]]-\frac1{12}[B,[A,B]]+\dots}
  \]

\end{frame}

%------------------------------------------------
\begin{frame}{Trotter Error}

  Using BCH

  \[
  e^{A+B}
  =
  e^A e^B e^{- \frac12[A,B] + \dots}
  \]

  Thus the error is determined by commutators.

\end{frame}

%------------------------------------------------
\begin{frame}{Leading Error}

  For first order Trotterization

  \[
  e^{(A+B)\Delta t}
  =
  e^{A\Delta t}e^{B\Delta t}
  +
  \mathcal{O}(\Delta t^2)
  \]

\end{frame}

%------------------------------------------------
\begin{frame}{Accuracy}

  Accuracy improves when

  \[
  \Delta t \rightarrow 0
  \]

  but

  \[
  n = \frac{1}{\Delta t}
  \]

  increases.

\end{frame}

%================================================
\section{Higher Order Suzuki Formulas}

%------------------------------------------------
\begin{frame}{Second Order Decomposition}

  Second order formula

  \[
  e^{(A+B)t}
  \approx
  e^{A t/2} e^{B t} e^{A t/2}
  \]

  Error

  \[
  \mathcal{O}(t^3)
  \]

\end{frame}

%------------------------------------------------
\begin{frame}{General Suzuki Formula}

  Higher order formulas exist

  \[
  S_{2k}(t)
  =
  S_{2k-2}(p_k t)^2
  S_{2k-2}((1-4p_k)t)
  S_{2k-2}(p_k t)^2
  \]

\end{frame}

%================================================
\section{Unitary Coupled Cluster}

%------------------------------------------------
\begin{frame}{UCC Ansatz}

  The unitary coupled cluster ansatz is

  \[
  |\Psi\rangle =
  e^{T-T^\dagger}|\Phi_0\rangle
  \]

\end{frame}

%------------------------------------------------
\begin{frame}{Cluster Operator}

  \[
  T_2 =
  \frac14
  \sum_{abij}
  t_{ij}^{ab}
  a_a^\dagger a_b^\dagger a_j a_i
  \]

\end{frame}

%------------------------------------------------
\begin{frame}{Generator}

  Define

  \[
  \tau = T-T^\dagger
  \]

  Thus

  \[
  U = e^{\tau}
  \]

\end{frame}

%------------------------------------------------
\begin{frame}{Operator Sum}

  \[
  \tau = \sum_k \theta_k G_k
  \]

  Each

  \[
  G_k
  \]

  represents a fermionic excitation.

\end{frame}

%================================================
\section{Jordan–Wigner Mapping}

%------------------------------------------------
\begin{frame}{Mapping Fermions to Qubits}

  \begin{align}
    a_p^\dagger &= \frac12(X_p-iY_p)\prod_{j<p}Z_j \\
    a_p &= \frac12(X_p+iY_p)\prod_{j<p}Z_j
  \end{align}

\end{frame}

%------------------------------------------------
\begin{frame}{Pauli Operators}

  Excitation operators become Pauli strings

  \[
  X_i X_j Y_k Y_l
  \]

\end{frame}

%================================================
\section{Quantum Circuit Blocks}

%------------------------------------------------
\begin{frame}{Circuit for Pauli Rotation}

  \[
  \begin{quantikz}
    \lstick{q_0} & \ctrl{1} & \qw & \qw & \qw \\
    \lstick{q_1} & \targ{} & \ctrl{1} & \qw & \qw \\
    \lstick{q_2} & \qw & \targ{} & \ctrl{1} & \qw \\
    \lstick{q_3} & \qw & \qw & \targ{} & \gate{R_z(\theta)}
  \end{quantikz}
  \]

\end{frame}

%================================================
\section{Circuit Depth Scaling}

%------------------------------------------------
\begin{frame}{Number of Excitations}

  For $N$ orbitals

  \[
  N_{exc} \sim N^4
  \]

\end{frame}

%------------------------------------------------
\begin{frame}{Total Depth}

  \[
  D \sim \frac{N^4}{\Delta t}
  \]

\end{frame}

%================================================
\section{Adaptive Ansatz}

%------------------------------------------------
\begin{frame}{ADAPT-UCC}

  ADAPT-UCC selects operators iteratively.

  Algorithm:

  \begin{itemize}
  \item compute gradient
  \item add operator with largest gradient
  \item reoptimize parameters
  \end{itemize}

\end{frame}

%================================================
\section{Parametric Matrix Models}

%------------------------------------------------
\begin{frame}{Motivation}

  UCC circuits scale poorly with system size.

  We introduce a reduced operator representation.

\end{frame}

%------------------------------------------------
\begin{frame}{PMM Ansatz}

  \[
  T(\theta)=\sum_k \theta_k M_k
  \]

\end{frame}

%------------------------------------------------
\begin{frame}{Reduced Circuit}

  \[
  U(\theta)
  =
  e^{T(\theta)-T^\dagger(\theta)}
  \]

\end{frame}

%------------------------------------------------
\begin{frame}{Scaling}

  Instead of

  \[
  N^4
  \]

  parameters we obtain

  \[
  N^2
  \]

\end{frame}

%================================================
\section{Summary}

%------------------------------------------------
\begin{frame}{Key Result}

  Circuit depth

  \[
  D \propto \frac{1}{\Delta t}
  \]

\end{frame}

%------------------------------------------------
\begin{frame}{Implications}

  For UCC

  \[
  D \sim \frac{N^4}{\Delta t}
  \]

  Strategies to reduce depth

  \begin{itemize}
  \item higher order Trotter
  \item adaptive ansätze
  \item parametric matrix models
  \end{itemize}

\end{frame}

\end{document}
